% =====================================================================
%  6. How fast does it converge?
% =====================================================================
\section{How fast does it converge?}

\begin{frame}{Order of convergence: GD is linear}
  \begin{columns}[T]
    \begin{column}{0.5\textwidth}
      \small
      Near a minimum GD behaves like $e_{k+1}\approx\rho\,e_k$ with $\rho = |1 - \alpha f''|$: the error shrinks by the
      \alert{same factor} every step. That is \textbf{linear} convergence (order $p = 1$).

      To cut the error by $10^{6}$ we need $\rho^k\le10^{-6}$:
      \[
        k \ge \frac{\ln 10^{6}}{\ln(1/\rho)} .
      \]
      \begin{itemize}\footnotesize
        \item $f = (x-3)^2$, $\alpha = 0.1$: $\rho = 0.8$, so $k = \ans{62}$ steps.
        \item running example, $\alpha = 0.1$: $\rho = 0.4$, so $k = \ans{16}$ steps.
      \end{itemize}
      Newton is \textbf{quadratic} ($p = 2$): correct digits roughly double each step.
    \end{column}
    \begin{column}{0.48\textwidth}
      \centering
      \includegraphics[width=\linewidth, height=0.56\textheight, keepaspectratio]{fig_gd_vs_newton}

      \footnotesize Log scale: linear falls steadily; quadratic falls faster
    \end{column}
  \end{columns}
\end{frame}

\begin{frame}{GD or Newton? Count the cost}
  \small
  \emph{When GD, and when Newton?} Count the work for $n$ parameters~\cite{nocedal2006}.

  \vspace{2pt}
  \centering\footnotesize
  \begin{tabular}{@{}l l l@{}}
    \toprule
    & Gradient descent & Newton's method \\
    \midrule
    information used & gradient ($n$ numbers) & gradient + Hessian ($n^2$ numbers) \\
    cost of one step & $O(n)$ after the gradient & $O(n^3)$ to solve $H\mathbf y = \g$ \\
    order of convergence & linear & quadratic \\
    learning rate & must be chosen & not needed \\
    on a concave stretch ($f'' < 0$) & still goes downhill & can head for a \textcolor{mRed}{maximum} \\
    best for & large problems (ML, $n\sim10^6$) & small problems needing high accuracy \\
    \bottomrule
  \end{tabular}

  \raggedright
  \vspace{3pt}
  \begin{keybox}[Why ML uses GD]
    \footnotesize With $n = 10^6$ a Hessian has $10^{12}$ entries and will not fit in memory; GD stores $n$ numbers.
    When $H$ is not positive definite, Levenberg--Marquardt~\cite{nocedal2006} moves towards a gradient step.
  \end{keybox}
\end{frame}

\statement{FIRST VS.\ SECOND ORDER}{First-order Taylor gives GD.\\ Second-order Taylor gives Newton.}{GD needs more steps; Newton does more work per step.}

\begin{frame}{Two assumptions: smooth, and strongly convex}
  \begin{columns}[T]
    \begin{column}{0.52\textwidth}
      \small
      \textbf{$L$-smooth}: the slope cannot change too fast,
      $\norm{\nabla J(\x) - \nabla J(\y)} \le L\norm{\x - \y}$. In 1-D: $|J''| \le L$.

      \vspace{3pt}
      \textbf{$\mu$-strongly convex}: the bowl is curved at least by $\mu$ everywhere. In 1-D: $J'' \ge \mu > 0$.

      \vspace{3pt}
      Together they squeeze $J$ between two parabolas through the current point. The
      \textbf{condition number}
      \[
        \kappa = L/\mu \ \ge 1
      \]
      measures how stretched the bowl is. It appears in every rate that follows~\cite{nesterov2018,boyd2004}.
    \end{column}
    \begin{column}{0.46\textwidth}
      \centering
      % Quadratic sandwich: mu-strong convexity (lower) and L-smoothness (upper) at x.
% f(y) = 0.5 y^2 + 0.6 log(cosh(1.5 y)), so f'' lies in [1, 1 + 1.35] = [1, 2.35]
\begin{tikzpicture}[scale=1.1, >={Stealth[length=5pt]},
  declare function={
    lc(\y)=ln(cosh(1.5*\y));
    f(\y)=0.5*\y*\y + 0.6*lc(\y);
    fp(\y)=\y + 0.9*tanh(1.5*\y);
    up(\y)=f(-0.6) + fp(-0.6)*(\y+0.6) + 0.5*2.35*(\y+0.6)*(\y+0.6);
    lo(\y)=f(-0.6) + fp(-0.6)*(\y+0.6) + 0.5*1.0*(\y+0.6)*(\y+0.6);}]
  % the corridor where f must live
  \fill[mBlue!8, domain=-1.7:1.1, samples=60]
    plot (\x, {up(\x)}) -- plot[domain=1.1:-1.7, samples=60] (\x, {lo(\x)}) -- cycle;
  \draw[->, mGrey] (-2.6,0) -- (2.4,0) node[right, font=\scriptsize] {$\y$};
  \draw[line width=1.1pt, mRed, domain=-1.7:1.1, samples=60] plot (\x, {up(\x)});
  \draw[line width=1.1pt, mBlue, domain=-2.2:2.0, samples=60] plot (\x, {lo(\x)});
  \draw[line width=1.9pt, mInk, domain=-2.2:2.0, samples=100] plot (\x, {f(\x)});
  \fill[mInk] (-0.6,{f(-0.6)}) circle (1.8pt);
  \draw[mGrey, densely dotted] (-0.6,0) node[below, font=\scriptsize, mInk] {$\x$} -- (-0.6,{f(-0.6)});
  \node[font=\scriptsize, mRed, anchor=south east, inner sep=3pt] at (1.1,{up(1.1)}) {$L$};
  \node[font=\scriptsize, mBlue, anchor=north west, inner sep=3pt] at (2.0,{lo(2.0)}) {$\mu$};
\end{tikzpicture}


      \scriptsize
      \textcolor{mRed}{\rule[2pt]{10pt}{1.1pt}} upper, curvature $L$ \quad
      \textcolor{mInk}{\rule[2pt]{10pt}{1.8pt}} $J$ \quad
      \textcolor{mBlue}{\rule[2pt]{10pt}{1.1pt}} lower, curvature $\mu$
    \end{column}
  \end{columns}
\end{frame}

\begin{frame}{The descent lemma}
  \begin{thmbox}[Descent lemma~\cite{nesterov2018}]
    If $J$ is $L$-smooth, then for all $\x, \y$:\quad
    $J(\y) \le J(\x) + \nabla J(\x)\T(\y - \x) + \tfrac L2\norm{\y - \x}^2$.
  \end{thmbox}
  \small
  \textbf{Proof idea.} Write $J(\y) - J(\x)$ as $\int_0^1 \nabla J(\x + t(\y - \x))\T(\y - \x)\,dt$ and bound the change in
  the gradient by $L\,t\norm{\y - \x}$. $\blacksquare$

  \vspace{4pt}
  Put in one GD step, $\y = \x_k - \alpha\g_k$:
  \[
    J(\x_{k+1}) \le J(\x_k) - \alpha\Bigl(1 - \frac{\alpha L}{2}\Bigr)\norm{\g_k}^2
    \;\overset{\alpha = 1/L}{=}\; J(\x_k) - \frac{1}{2L}\norm{\g_k}^2 .
  \]
  \begin{keybox}[What it tells us]
    \small Any $\alpha < 2/L$ is \emph{guaranteed} to lower $J$, the same window as in Section 5, now for any smooth $J$.
    Each step removes at least $\norm{\g_k}^2/2L$, and every rate below is built on that drop~\cite{bubeck2015}.
  \end{keybox}
\end{frame}

\begin{frame}{Strongly convex: the error shrinks by a fixed factor}
  \begin{thmbox}[Theorem ($L$-smooth, $\mu$-strongly convex, $\alpha = 1/L$)]
    \[ \norm{\x_k - \x^\star}^2 \le \Bigl(1 - \frac1\kappa\Bigr)^{k}\norm{\x_0 - \x^\star}^2 . \]
  \end{thmbox}
  \small\textbf{Proof.} Let $\e = \x_k - \x^\star$ and $\Delta = J(\x_k) - J^\star$. Expand one step:
  $\norm{\x_{k+1} - \x^\star}^2 = \norm{\e}^2 - \tfrac2L\g_k\T\e + \tfrac1{L^2}\norm{\g_k}^2$. Two facts:
  \begin{itemize}\footnotesize
    \item strong convexity gives $\g_k\T\e \ge \Delta + \tfrac\mu2\norm{\e}^2$;
    \item the descent lemma gives $J^\star \le J(\x_k) - \tfrac1{2L}\norm{\g_k}^2$, so $\norm{\g_k}^2 \le 2L\Delta$.
  \end{itemize}
  \[
    \norm{\x_{k+1} - \x^\star}^2 \le \norm{\e}^2 - \tfrac2L\Delta - \tfrac\mu L\norm{\e}^2 + \tfrac2L\Delta
    = \Bigl(1 - \tfrac1\kappa\Bigr)\norm{\e}^2 . \quad\blacksquare
  \]
  \footnotesize This is the \emph{linear} order from the start of this section, now proved for every smooth, strongly convex function. It needs $O(\kappa\log\tfrac1\varepsilon)$ steps.
\end{frame}

\begin{frame}{Convergence guarantees}
  \begin{columns}[T]
    \begin{column}{0.54\textwidth}
      \footnotesize
      \begin{tabular}{@{}l l l@{}}
        \toprule
        Assumption & Guarantee ($\alpha = \frac1L$) & Steps to $\varepsilon$ \\
        \midrule
        smooth only & $\min_k\norm{\g_k}^2 \le \frac{2L\Delta_0}{K}$ & $O(\varepsilon^{-2})$ \\
        convex & $J - J^\star \le \frac{LR^2}{2k}$ & $O(\varepsilon^{-1})$ \\
        strongly convex & $\norm{\e_k}^2 \le (1 - \frac1\kappa)^k R^2$ & $O(\kappa\log\frac1\varepsilon)$ \\
        PL~\cite{karimi2016} & $J - J^\star \le (1 - \frac1\kappa)^k\Delta_0$ & $O(\kappa\log\frac1\varepsilon)$ \\
        \bottomrule
      \end{tabular}

      \vspace{4pt}
      \begin{itemize}\footnotesize
        \item Smooth only (deep learning): GD reaches a point with a small gradient, possibly a saddle.
        \item PL, $\tfrac12\norm{\nabla J}^2 \ge \mu(J - J^\star)$, needs no convexity and still gives a linear rate.
        \item $R = \norm{\x_0 - \x^\star}$, $\Delta_0 = J(\x_0) - J^\star$.
      \end{itemize}
    \end{column}
    \begin{column}{0.44\textwidth}
      \centering
      \includegraphics[width=\linewidth, height=0.62\textheight, keepaspectratio]{fig_rate_bounds}
    \end{column}
  \end{columns}
\end{frame}
